\section{Short-shift moments after weighted transfer}
\label{sec:short-shift-moments}

For the remainder of the positive theory, fix numbers
$0<\eps<13/15$ and $0<\eta<1$.  Thus
\begin{equation}
 \frac2{15}+\eps<1,
 \label{eq:nonempty-short-range}
\end{equation}
and assume
\begin{equation}
 X^{2/15+\eps}\leq H\leq X,
 \qquad
 1\leq D\leq X^{1-\eta}.
 \label{eq:main-short-range}
\end{equation}
The support of $G$ is contained in $(0,\infty)$, so the average samples a
growing interval $h\asymp_GH$; it does not contain any fixed shift once
$X$ is large.

\subsection{The first-moment model}

\begin{theorem}[Short-shift first moment]
\label{thm:short-first-moment}
For every $A>0$, uniformly for \eqref{eq:main-short-range} and all
$t\in\R$,
\begin{equation}
 \frac1H\sum_hG(h/H)\cT_{h,D}(X;t)
 =I_0(G)\cM_D(X;t)
 +O_{A,\eps,W,G}\bigl(X(\log X)^{-A}\bigr).
 \label{eq:short-first-moment}
\end{equation}
Equivalently, the normalized smooth average of $\cR_{h,D}$ has the
displayed error.
\end{theorem}

\begin{proof}
Group the residual average by $r=mn$.  Equations
\eqref{eq:intro-centered-transform} and
\eqref{eq:product-signed-weight} give
\begin{equation}
 \frac1H\sum_hG(h/H)\cR_{h,D}(X;t)
 =\sum_rB_{D,t}(r)
 \left\{\frac1H\sum_hG(h/H)(\Lambda(r+h)-1)\right\}.
 \label{eq:first-moment-grouped}
\end{equation}
On the complement of the exceptional set in
\cref{lem:smooth-almost-all-moments}, the braces are
$O_A((\log X)^{-A})$.  The pointwise estimate
\eqref{eq:B-pointwise} and the classical mean order of $\tau$ give
\[
 \sum_r|B_{D,t}(r)|\ll_WX\log(2X).
\]
Request two additional logarithmic powers.

On the exceptional set, the braces are $O_G(\log(2X))$ by the elementary
upper-bound sieve for primes in an interval; the completely trivial
$O_G((\log X)^2)$ would also suffice.  Apply
 \eqref{eq:B-exception-transfer} with a larger value of $A$.  Both parts
 are uniform in $t$ and show that the normalized average of
 $\cR_{h,D}$ is $O_A(X(\log X)^{-A})$.  Finally,
 \[
  \frac1H\sum_hG(h/H)=I_0(G)+O_G(H^{-1})
 \]
 and the elementary bound $|\cM_D(X;t)|\ll_WX\log(2X)$ shows that the
 Riemann-sum error is absorbed by the stated remainder, since
 $H\geq X^{2/15+\eps}$.  This proves the displayed formula for
 $\cT_{h,D}$.
\end{proof}

\subsection{The transferred target square}

\begin{proposition}[Weighted target-square transfer]
\label{prop:weighted-target-square}
Uniformly in \eqref{eq:main-short-range},
\begin{equation}
 \begin{aligned}
 &\sum_{r\geq1}A_D(r)
 \sum_hG(h/H)(\Lambda(r+h)-1)^2\\
 &\qquad={}
 HI_0(G)\log X\,\cS_D(X)
 +O_{\eps,W,G}\bigl(HX(\log X)^2\bigr).
 \end{aligned}
 \label{eq:weighted-target-square}
\end{equation}
The same formula holds with the target square replaced by
$\Lambda(r+h)^2$.
\end{proposition}

\begin{proof}
For $r$ outside the exceptional set, use
\eqref{eq:smooth-eta-square-moment}; the total of its $O(H)$ error is
$O(H\cS_D(X))=O(HX(\log X)^2)$ by
\cref{lem:source-diagonal}.

On the exceptional set, the whole inner target-square sum is at most
$O_G(H(\log(2X))^2)$.  The same bound covers the main term being
removed there.  Proposition~\ref{prop:weighted-exception-transfer}, with
an arbitrarily large logarithmic saving, makes the exceptional
contribution $O(HX(\log X)^{-10})$, say.  This proves the centered
formula.  Equation \eqref{eq:smooth-Lambda-square-moment} proves the raw
version in exactly the same way.
\end{proof}

The proposition is the point where the almost-all theorem becomes a
complete asymptotic input: no product center is discarded from the final
sum.

\subsection{A uniform two-target upper sieve}

For $q\ne0$, put
\begin{equation}
 \mathfrak S^+(q)=
 \prod_{\substack{p\mid q\\p>2}}\frac{p-1}{p-2}.
 \label{eq:upper-singular-factor}
\end{equation}

\begin{lemma}[Two translated targets in a short shift window]
\label{lem:short-two-target-sieve}
Let $r\ne r'$ lie in $[\alpha X,\beta X]$.  Uniformly for
$H\geq X^{2/15+\eps}$ and $H\leq X$,
\begin{equation}
 \begin{aligned}
 &\sum_hG(h/H)
 \left|(\Lambda(r+h)-1)(\Lambda(r'+h)-1)\right|\\
 &\qquad\ll_{\eps,W,G}H\mathfrak S^+(r-r')
 \ll_{\eps,W,G}H\log\log(3X).
 \end{aligned}
 \label{eq:short-two-target-sieve}
\end{equation}
The first bound remains valid if one or both centered factors are
replaced by $\Lambda$.
\end{lemma}

\begin{proof}
The Selberg upper-bound sieve for the two linear forms $h+r$ and $h+r'$
gives
\[
 \#\{h\asymp_GH:h+r,\ h+r'\text{ are prime}\}
 \ll_G\frac{H}{(\log H)^2}\mathfrak S^+(r-r').
\]
Restoring the von Mangoldt weights costs $(\log X)^2$.  Since
$H\geq X^{2/15+\eps}$, the ratio $\log X/\log H$ is bounded in terms of
$\eps$, and the prime-prime contribution has the asserted size.  The
constant and one-prime pieces follow from the corresponding one-linear-
form upper sieve.

If one target is a proper prime power, the local spacing argument used in
\cref{lem:smooth-almost-all-moments} gives only a polylogarithmic term
plus $O(HX^{-1/2}(\log X)^2)$, both $o_\eps(H)$.  Finally the standard
maximal-order estimate
\[
 \mathfrak S^+(q)\ll\log\log(3|q|)
\]
follows by comparing with the product over the smallest primes.  Here
$0<|r-r'|\ll_WX$.  See \citet[Chapter~6]{HalberstamRichert1974} for the
upper sieve.
\end{proof}
